% FONTI E NOTE DI LAVORO DEL CAPITOLO DEI RISULTATI
% Questo file non viene compilato: raccoglie, per ogni sottosezione della
% versione dell'11 ottobre 2026 (commit 48a4c9d), i commenti con l'origine
% dei valori (file sul server, run, data di lettura), i punti da verificare
% e i testi tolti dal capitolo. I titoli sono quelli di quella versione.
% Contiene percorsi del server e nomi di run: da togliere se il repository
% diventa pubblico.

% ======================================================================
% Note iniziali
% ======================================================================
% Questa scaletta non contiene ancora risultati numerici dei nuovi confronti.
% Integrare soltanto misure riferite a esecuzioni identificate e confrontabili.

% ======================================================================
% Accuratezza del forecasting
% ======================================================================
% Il backbone deterministico non ha una run propria su PVGIS con il protocollo
% delle varianti SDE (60 epoche): esiste sui dati reali, esclusi dalla tesi, e su
% PVGIS come membro del Deep Ensemble (5 epoche, MAE circa 20 W). Non presentarlo
% come riferimento alla pari delle varianti SDE. Una run dedicata resta
% un'opzione aperta.

% ======================================================================
% Dropout Monte Carlo e Deep Ensemble
% ======================================================================
% NUMERI LETTI SU W&B il 10 ottobre 2026 (progetto PhysiQ-PV, run del 10 giugno
% 2026): enhanced-daytime-check-seed1 (id x4eic20l), MAE 20.1 W; aggregazione
% pvgis_deep_ensemble_vfry1hgx (id 29yahjvt), cinque membri, 10037664 campioni,
% MAE 19.19 W, RMSE 45.06 W. DA FARE: RMSE del dropout Monte Carlo e metriche
% sulle sole ore diurne, da rileggere su W&B.

% ======================================================================
% Variante SDE U-Net
% ======================================================================
% NUMERI LETTI SUL SERVER l'11 ottobre 2026, file metrics_global.csv:
%   pvgis_stgnn_normal_only_monaco_ep60_seed42: MAE 18.672, RMSE 46.090
%   pvgis_stgnn_normal_only_sigma1_ep60_seed42: MAE 18.619, RMSE 46.125,
%     copertura 0.126, deviazione standard media 2.876 W. Seconda esecuzione
%     SDE U-Net: DA VERIFICARE che cosa cambia ("sigma1") prima di citarla.
%   pvgis_stgnn_paper_faithful_gaussian_ep60_seed1: MAE 20.638, RMSE 49.162
%   ..._detector_mtgflow_ep60_seed1: MAE 20.652, RMSE 49.112
%   ..._detector_mtgflow_ep60_h1-2-3-4-5-6_direct_seed1: MAE 32.272, RMSE 69.045
%     su 60191514 previsioni (10031919 per orizzonte).
% Tutti in \provv: le run hanno training 2016--2017, validation 2018 e il
% riferimento a cielo sereno con inclinazione di 30 gradi.
% Ore diurne e notturne (metrics_daytime.csv, righe daytime e nighttime):
%   SDE U-Net: diurne MAE 38.064, RMSE 64.981; notturne MAE 1.269, RMSE 15.531
%   SDE-Net:   diurne MAE 42.795, RMSE 69.611; notturne MAE 0.753, RMSE 15.403
% Il RMSE notturno di circa 15 W contro un MAE di circa 1 W indica poche ore
% "notturne" con errori grandi (irradianza sotto 10 W/m2 ma produzione non
% nulla): DA GUARDARE prima di commentarlo; esiste pvgis_data_quality_issues.csv.
% Righe high_daytime (1186852), peak_daytime (474719) ed extreme_peak_daytime
% (237361): soglie non note, DA LEGGERE nel codice prima di usarle. In entrambe
% le varianti la produzione alta è sottostimata (residuo medio fra -29 e -44 W,
% quota di sottostime fra 0.74 e 0.94).

% ======================================================================
% Variante SDE-Net e previsione diretta multiorizzonte
% ======================================================================
% Fonte: multihorizon_anomaly_group_metrics.csv della run a sei orizzonti (letto
% l'11 ottobre 2026). Filtro diurno: multihorizon_posthoc_metadata.json non lo
% riporta. Indizi concordi: 4761227 righe per orizzonte sono il 47.5% delle
% 10031919 previsioni (47.3% di ore diurne nel modello a un'ora), e l'analisi di
% sensibilità, che dichiara la soglia diurna di 10 W/m2, dà a un'ora e k = 1.5
% valori quasi uguali (normali 44.96, segnalate 51.01 W, su 4724704 righe).
% DA CONFERMARE nel notebook che ha prodotto il file.
% Non confondere t+1 e t+6 dello stesso modello diretto con risultati di modelli
% addestrati separatamente.
% Sottosezione tolta l'11 ottobre 2026: nessuna run senza GAT sul server (vedi
% il commento nella Metodologia, sezione sul grafo geografico).

% ======================================================================
% Ablazione della GAT
% ======================================================================
%
% Appartiene alla variante SDE-Net. Confrontare configurazioni e campioni
% coerenti. L'ablazione rimuove il message passing, ma conserva la diffusione
% globale condivisa fra località.

% ======================================================================
% Intervalli del dropout Monte Carlo e del Deep Ensemble
% ======================================================================
% NUMERI LETTI SU W&B il 10 ottobre 2026: copertura 0.028 (dropout Monte Carlo,
% un modello, 20 passaggi) e 0.173 (Deep Ensemble, cinque membri); deviazione
% standard media del Deep Ensemble 5.78 W.
% DA FARE: ampiezza media, NMPIL e CLC ricalcolato con eta = 9 (su W&B il CLC è
% calcolato con eta = 10); deviazione standard del dropout Monte Carlo. Le
% metriche sulle sole osservazioni segnalate NON vanno prese da queste run: le
% loro etichette usano una climatologia diversa (vedi Capitolo 5).

% ======================================================================
% Intervalli delle due varianti SDE
% ======================================================================
% Fonte: posthoc_by_horizon/t_plus_1 e t_plus_6 della run a sei orizzonti
% (letti l'11 ottobre 2026).
% uncertainty_components.csv, ore diurne (4747545): deviazione standard media
% epistemica, aleatoria, totale; quota epistemica della varianza.
%   t+1 tutte      0.04429,  62.7095,  62.7096, 6.68e-06
%   t+1 normali    0.04288,  62.3454,  62.3455, 6.28e-06  (4434310)
%   t+1 segnalate  0.06430,  67.8637,  67.8640, 1.24e-05  (313235)
%   t+6 tutte      0.03106, 102.7968, 102.7968, 1.01e-06
%   t+6 normali    0.02948, 102.2012, 102.2012, 8.90e-07
%   t+6 segnalate  0.05343, 111.2284, 111.2284, 2.77e-06
% daytime_bin_summary.csv: numero, MAE, RMSE, PICP, deviazione standard media,
% ampiezza media, NMPIL.
%   t+1  0-20   1934608, 39.914,  66.973, 0.9343,  44.917, 176.070, 0.1974
%   t+1 20-40    856630, 52.211,  71.526, 0.9485,  71.202, 279.105, 0.3129
%   t+1 40-60    736468, 51.497,  73.936, 0.9586,  78.407, 307.349, 0.3446
%   t+1 60-80    669745, 48.930,  78.807, 0.9691,  77.575, 304.087, 0.3410
%   t+1 80-100   550094, 44.578,  74.743, 0.9884,  72.948, 285.951, 0.3206
%   t+6  0-20   71.826, 119.323, 0.9073,  62.786, 246.116, 0.2760
%   t+6 20-40   75.655, 102.070, 0.9676, 100.050, 392.187, 0.4397
%   t+6 40-60   84.072, 100.565, 0.9885, 125.060, 490.227, 0.5497
%   t+6 60-80   92.056, 111.779, 0.9952, 144.923, 568.087, 0.6370
%   t+6 80-100 117.075, 142.628, 0.9949, 166.693, 653.423, 0.7327
% Picco di riferimento per le fasce: 712.95 W (analysis_metadata.json).
% La separazione delle componenti è disponibile solo per la run a sei orizzonti;
% per quella a un'ora andrebbe letta da predictions.csv.
% NUMERI LETTI SUL SERVER l'11 ottobre 2026, metrics_daytime.csv, colonne
% picp_pi, mpiw_pi, nmpil_pi, clc_pi (intervallo principale), righe all e daytime:
%   SDE U-Net: all 0.10805, 8.184, 0.009161, 17.908; daytime 0.22724, 16.817,
%     0.018827, 12.602; nighttime 0.00109, 0.435
%   SDE-Net:   all 0.97654, 108.200, 0.121129, 0.21652; daytime 0.95311,
%     224.562, 0.251393, 0.49585; nighttime 0.99758, 3.774
% Il CLC dei file è calcolato con eta = 9 (verificato: 0.009161 * (1 +
% exp(9 * 0.84195)) = 17.91). R = 893.27 W.
% ATTENZIONE: per SDE U-Net la colonna coverage_95 di metrics_global.csv e di
% metrics_by_anomaly_label.csv (0.122 su tutte le ore) è quella della banda
% gaussiana media +/- 1.96 deviazioni standard (picp_gaussian), non quella dei
% quantili empirici (picp_pi, 0.108). Per SDE-Net le due coincidono (0.9765).
% Deviazione standard (metrics_global.csv: media, mediana, percentile 90):
%   SDE U-Net 2.386, 0.333, 6.226; SDE-Net un'ora 27.603, 0.786, 87.759;
%   SDE-Net sei orizzonti 41.975, 0.841, 131.392 (copertura 0.9770).
% Costruzione degli intervalli letta nel codice (training/uncertainty.py dei due
% branch): quantili empirici delle traiettorie per SDE U-Net; quantili della
% miscela gaussiana per SDE-Net, con deviazione standard totale
% sqrt(varianza epistemica + varianza aleatoria).
% Componenti epistemica e aleatoria: vedi il commento dopo la Tabella delle
% fasce di produzione, più sotto.
% Per la variante SDE-Net distinguere quantili della miscela gaussiana e banda
% gaussiana approssimata; riportare separatamente le componenti della varianza.
% ESITO DA SOSTENERE CON I NUMERI (spostato qui dallo Stato dell'arte): la testa
% gaussiana avrebbe dato intervalli utilizzabili, la banda epistemica (dispersione
% fra le traiettorie) sarebbe troppo stretta, in modo più evidente nella variante
% SDE U-Net, che non ha componente aleatoria. Ordine di grandezza ricordato
% dall'autore: la banda epistemica conteneva circa il 10--15% delle osservazioni.
% Ritrovare nei file di valutazione copertura, ampiezza, livello nominale e
% variante prima di scriverlo. Riprendere il limite anche nelle Conclusioni.

% ======================================================================
% Score e distribuzione spazio-temporale delle anomalie
% ======================================================================
% TESTO SPOSTATO DAL CAPITOLO 5 (11 ottobre 2026). Descrive come è calcolata
% l'analisi: da riusare, accorciato, in apertura della sottosezione e soltanto
% insieme al risultato corrispondente.
% Le previsioni sono associate alle decisioni dei detector tramite località e timestamp target. Nel caso multiorizzonte, la chiave completa della previsione include anche l'orizzonte, mentre l'etichetta del detector rimane unica per località e timestamp. La stessa osservazione target può quindi essere associata a previsioni emesse in ore diverse.
%
% Il protocollo STGAN prevede un'esclusione opzionale degli azzeramenti solari regionali isolati e del passo immediatamente successivo, seguito dal ricalcolo del top-K sulle osservazioni ammissibili. Il percorso NetCDF controlla più variabili solari e la produzione; quello basato sui dati preparati controlla la sola irradianza. La sorgente utilizzata deve essere indicata, poiché i due criteri non sono identici.
%
% La serie locale mostra tempo e anomaly score, con le decisioni evidenziate e interruzioni nei periodi senza score. La rappresentazione regionale mostra invece, per ogni ora, la percentuale di siti anomali rispetto ai siti con score valido; una matrice sito--giorno mostra la percentuale di ore anomale rispetto alle ore valide di quel sito e giorno. Queste quote aggregano le decisioni puntuali senza stimare una nuova soglia regionale. I periodi senza osservazioni non vengono convertiti in normalità.
% Fonte: pvgis_stgan_cnn/convgru_reference/seed_20/entity_anomaly_scores.csv
% (10 ottobre 2026), colonna contribution_fraction, 100653 righe per variabile.
% Dal riepilogo per località (summary_by_seed.csv, visto in parte): la quota di
% osservazioni segnalate per località varia circa da 0.4% a 3.0%.
% Inserire score locali e quote regionali indicando copertura e denominatori.
% Distinguere score aggregati del detector e diagnostiche per variabile.

% ======================================================================
% Soglie IQR, top-K e confronto fra grafo e griglia
% ======================================================================
% Fonte: reference_detector_overlap.csv (letto l'11 ottobre 2026), soglie di
% riferimento k = 1.5 e 1%: n_common_daytime 4724704, n_mtgflow_rare 312044,
% n_stgan_rare 94532, n_both_rare 6829, jaccard 0.017083,
% stgan_captured_by_mtgflow 0.07224, mtgflow_captured_by_stgan 0.02188,
% daily_pearson 0.13678, daily_spearman 0.18259.
% Atteso per caso: 312044 * 94532 / 4724704 = 6243; Jaccard 6243 / (312044 +
% 94532 - 6243) = 0.0156.
% STGAN = versione su grafo: confermato dal protocollo dell'analisi dei quattro
% casi, che usa pvgis_stgan/paper_reference/seed_20 (stessa serie di analisi).
% DA VERIFICARE nel notebook: quale grandezza giornaliera è correlata (file
% reference_detector_daily_overlap.csv) e che STGAN sia la versione su grafo.
% reference_decision_metrics.csv, errore assoluto a un'ora (media; primo
% quartile, mediana, terzo quartile): MTGFlow normali 44.96; 10.76, 27.35, 58.18;
% segnalate 51.01; 12.22, 30.12, 66.10. STGAN normali 45.04; 10.76, 27.32, 58.29;
% segnalate 61.19; 17.73, 37.99, 76.74. A sei ore: MTGFlow 82.12 (22.97, 59.58,
% 113.54) e 89.00 (24.82, 63.32, 123.05); STGAN 82.35 (22.99, 59.68, 113.85) e
% 93.80 (28.48, 66.91, 129.16).
% Fonte: pvgis_mtgflow/downstream_dense/summary_aggregate.csv, colonna
% anomaly_rate_mean, 1149 località (letto l'11 ottobre 2026): minimo 0.01333,
% primo quartile 0.04747, mediana 0.06045, terzo quartile 0.07390, massimo
% 0.13677; 8701 osservazioni con score per località (8760 meno le prime 59 ore,
% finestra di 60). Quota complessiva: 613164 segnalate su 9997449 con score.
% Per STGAN ConvGRU, dal riepilogo per località visto in parte, la quota per
% località varia circa da 0.4% a 3.0%: DA RILEGGERE prima di scriverlo.
% TESTO SPOSTATO DAL CAPITOLO 5 (11 ottobre 2026). Descrive come è calcolata
% l'analisi: da riusare, accorciato, in apertura della sottosezione e soltanto
% insieme al risultato corrispondente.
% Le mappe di frequenza associano a ciascun sito la quota di ore anomale sulle ore con score valido. La copertura è distinta da questa quota e rapporta le ore effettivamente valutate a quelle ammissibili. Un sito privo di score non viene rappresentato come privo di anomalie. Nel confronto diurno, il filtro dell'irradianza maggiore di $10\,\mathrm{W\,m^{-2}}$ viene applicato dopo le decisioni del detector: il ranking STGAN non viene ricalcolato sulle sole ore diurne.
%
% Per MTGFlow, il protocollo di base conserva le soglie per sito stimate dalle statistiche di training. Una variante esplorativa assegna invece un budget globale di anomalie ordinando la coordinata firmata
% \begin{equation}
% c_{nt}=\frac{s_{nt}-\theta_n}{\mathrm{IQR}_{n,\mathrm{train}}},
% \end{equation}
% dove $s_{nt}$ è lo score e $\theta_n$ la soglia salvata del sito. Non viene applicato un troncamento a zero: anche i valori sotto soglia mantengono il proprio ordinamento. La selezione usa un numero esatto di osservazioni determinato dalla percentuale richiesta, prima dell'eventuale filtro diurno. Questa variante deve essere distinta sia dalle decisioni native MTGFlow sia dalla coordinata centrata in $Q_3$ usata nell'analisi di sensibilità.
% Specificare statistiche IQR per sito e popolazione del ranking globale.
% Per STGAN su grafo e ConvGRU riportare accordo e Jaccard sulle coordinate
% comuni, separando patch complete/incomplete. Non interpretarli come
% accuratezza del detector.
% Non mescolare flag salvati con top-K ricalcolato dopo le esclusioni.
% Distinguere le soglie native MTGFlow dalla variante esplorativa a budget
% globale sulla coordinata firmata (score - soglia salvata) / IQR del training.
% Nelle mappe riportare frequenza e copertura con i rispettivi denominatori;
% il filtro diurno successivo alle decisioni non ricalcola il ranking.

% ======================================================================
% Stabilità delle decisioni al variare del modello
% ======================================================================
% NUMERI VERIFICATI (server, 10 ottobre 2026): confronto di ciascuna run ConvGRU
% con la versione su grafo (pvgis_stgan/paper_reference), file
% seed_20/baseline_comparison/detector_comparison_by_boundary.csv, decisioni
% salvate top 1% globale senza filtri né ricalcolo del ranking.
% Osservazioni: 10065240 (1149 località x 8760 ore); segnalate da ciascun
% detector: 100653. Con due detector indipendenti all'1% l'intersezione attesa
% per caso è circa 1007 osservazioni (Jaccard circa 0.005).
%   kernel 1: 9956 in comune, Jaccard 0.0520 (patch complete 0.0537, incomplete 0.0463)
%   kernel 3: 6267 in comune, Jaccard 0.0321 (patch complete 0.0322, incomplete 0.0319)
%   kernel 5: 7325 in comune, Jaccard 0.0378 (patch complete 0.0388, incomplete 0.0342)
% Quota segnalata sulle patch incomplete (232 località) e complete (917):
%   grafo 1.12% e 0.97%; kernel 1 1.11% e 0.97%; kernel 3 1.27% e 0.93%;
%   kernel 5 1.07% e 0.98%.
% Loss per epoca (seed_20/training_history.csv), sei epoche, 140803056 campioni
% per epoca (1149 località x 122544 istanti: nessun sottocampionamento):
%   kernel 3: G 2.99 3.18 3.23 3.32 3.46 3.47; D 0.629 0.622 0.618 0.612 0.603 0.601
%   kernel 1: G 3.12 3.21 3.74 3.42 3.49 3.38; D 0.597 0.596 0.579 0.595 0.595 0.601
%   kernel 5: G 2.96 2.95 3.17 3.58 3.33 3.14; D 0.633 0.632 0.623 0.614 0.613 0.626
% Parametri con kernel 3 (model_results_summary.txt): generatore 140931,
% discriminatore 114593. Da leggere per kernel 1 e 5.
% ConvGRU kernel 3, statistiche dello score (somma di due componenti in [0,1]):
%   patch complete: media 0.429, mediana 0.489, quantile 0.95 0.553, 74776 segnalate (0.93%)
%   patch incomplete: media 0.430, mediana 0.492, quantile 0.95 0.563, 25877 segnalate (1.27%)
% Training 51067 s, scoring 278 s, memoria GPU di picco 485 MB.
% Le due componenti dello score non sono state salvate: il peso di ciascuna sullo
% score non è verificabile senza rifare lo scoring. La mediana vicina a 0.49 è
% compatibile con una differenza D(reale) - D(generato) vicina a zero per
% l'osservazione tipica (deduzione, non misura).
% entity_anomaly_scores.csv riporta, per le osservazioni segnalate, la quota
% dell'errore di ricostruzione dovuta a ciascuna variabile.
% Tempo per epoca: circa 8280 s (kernel 3), 7950 s (kernel 1), 3430 s (kernel 5
% con caricamento dei dati ottimizzato). La versione su grafo non ha salvato la
% storia delle loss. La loss di D resta vicina a 0.693, valore di un
% discriminatore che non distingue reale e generato.
% Limiti da dichiarare: un solo seme (20) per ogni configurazione, quindi la
% variabilità fra semi non è misurata e la differenza non si può attribuire al
% solo operatore; manca il confronto diretto fra kernel (solo kernel contro grafo).
% Inserire soltanto confronti fra run completate e identificate per k=1,3,5.
% Tenere distinta l'ampiezza P della patch; il default P=3 non fornisce nuovi
% siti osservati quando k passa a 5. Riportare i parametri di G e D per variante.
% Il caso k=1 elimina il mixing convoluzionale, ma conserva aggregazioni e
% loss sulla patch: non presentarlo come assenza completa di contesto spaziale.
% Confrontare campioni, maschere, finestre, seed e budget di training coerenti.
% Non dedurre la superiorità di un kernel dalla sola disponibilità del codice.

% ======================================================================
% MAE e RMSE sui campioni normali e anomali
% ======================================================================
% Fonte: metrics_daytime.csv, righe normal_daytime e rare_extreme_daytime
% (l'11 ottobre 2026). Rapporti: 43.169/37.339 = 1.156; 50.910/41.642 = 1.223.
% Quote diurne: 590695/1059283 = 55.8%; 4156850/8978381 = 46.3%.
% DA FARE: stesso confronto per fascia di produzione
% (residual_bias_and_bin_metrics.csv) e per tipo di etichetta.
% NUMERI LETTI SUL SERVER l'11 ottobre 2026, file metrics_by_anomaly_label.csv
% delle run normal_only_monaco_ep60_seed42 (SDE U-Net),
% paper_faithful_gaussian_ep60_seed1 (SDE-Net, etichette climatologiche),
% ..._detector_mtgflow_ep60_seed1 e ..._h1-2-3-4-5-6_direct_seed1 (etichette
% MTGFlow). Righe "group:" = etichetta locale dell'osservazione; "label:" = tipo.
% Nel codice l'etichetta solar_potential copre irradianza e produzione
% (pvgis_anomaly_scores.py, SOLAR_VARIABLES).
% Non usate qui le righe "event:" (eventi regionali, criterio del filtro sulle
% finestre normali). Per SDE U-Net: ore in un evento regionale 785916 (7.8%),
% MAE 41.09 e RMSE 85.73 contro 16.77 e 40.99 nelle altre; copertura 0.120 e
% 0.122. Con le etichette MTGFlow, un'ora: evento 346998, MAE 28.02, RMSE 58.60.
% Rapporti: 27.61/17.62 = 1.57; 31.65/19.34 = 1.64; 26.01/20.30 = 1.28;
% 38.58/31.86 = 1.21; 98.51/17.62 = 5.6; 112.67/19.34 = 5.8.
% Riportare conteggi, metriche aggregate e distribuzioni degli errori assoluti
% per fascia di produzione. I boxplot non descrivono RMSE per campione.

% ======================================================================
% Confronto a una e sei ore e quattro casi input--target
% ======================================================================
% Differenze: 51.976 - 45.186 = 6.79; 88.910 - 81.836 = 7.07.
% Fonte: stgan_input_target_cases_t1_t6 (letto l'11 ottobre 2026).
% Protocollo: sola rielaborazione; previsioni della run a sei orizzonti; score
% pvgis_stgan/paper_reference/seed_20 (versione su grafo); top 1% ricalcolato
% dopo il filtro di qualità, soglia dello score 0.65091; finestra di 24 ore
% della stessa località che termina all'istante di emissione; almeno un passo
% anomalo rende anomala la finestra; soglia diurna 10 W/m2.
% Righe diurne caricate 9495090 (due orizzonti), classificate 9217715, escluse
% 277375. A un'ora: incluse 4611053 (97.125%), finestra incompleta 126230,
% finestra incompleta e target senza score 10262. A sei ore: incluse 4606662
% (97.0325%), finestra incompleta 130621, entrambe 2298, target senza score 7964.
% Casi (numero, quota, MAE, RMSE):
%   t+1 normale-normale 3570325, 77.430%, 42.902,  67.197
%   t+1 normale-anomalo   66845,  1.450%, 59.421,  93.254
%   t+1 anomalo-normale  948460, 20.569%, 50.326,  76.465
%   t+1 anomalo-anomalo   25423,  0.551%, 62.205,  97.524
%   t+6 normale-normale 3567175, 77.435%, 79.675, 111.993
%   t+6 normale-anomalo   68398,  1.485%, 92.733, 129.407
%   t+6 anomalo-normale  947430, 20.567%, 90.199, 125.021
%   t+6 anomalo-anomalo   23659,  0.514%, 95.190, 132.222
% Rapporti rispetto a normale-normale: t+1 1.173, 1.385, 1.450; t+6 1.132,
% 1.164, 1.195.
% DA FARE: stessi casi per fascia di produzione (input_target_bin_metrics.csv),
% per controllare che l'effetto dell'ingresso anomalo non dipenda dalla diversa
% ripartizione fra le fasce.
% TESTO SPOSTATO DAL CAPITOLO 5 (11 ottobre 2026). Descrive come è calcolata
% l'analisi: da riusare, accorciato, in apertura della sottosezione e soltanto
% insieme al risultato corrispondente.
% Oltre al raggruppamento secondo l'etichetta del target, l'analisi considera quattro casi: input normale e target normale, input normale e target anomalo, input anomalo e target normale, input anomalo e target anomalo. La normalità del target riguarda l'osservazione futura classificata dal detector, non una classificazione della previsione numerica del modello.
%
% Con una sequenza di 24 ore e ultimo input al tempo $t$, la finestra considerata è $[t-23\,\mathrm{h},t]$, estremi inclusi. Nel criterio predefinito, almeno un timestamp anomalo rende anomala la finestra. Sono considerate tutte le ore disponibili, anche notturne, della stessa località del target.
%
% La classificazione descrive la sequenza esplicita del sito centrale: non comprende automaticamente le altre località del grafo né le ore aggiuntive necessarie alle feature derivate. Le finestre incomplete e i target privi di score sono esclusi con conteggi separati. Il numero di righe indica previsioni località--timestamp per orizzonte, non giorni statisticamente indipendenti.
% Inserire i quattro casi con copertura ed esclusioni della finestra di input.
% Il target normale/anomalo è l'osservazione etichettata dal detector:
% questo confronto non è una matrice di confusione del previsore.

% ======================================================================
% Due episodi del 2019
% ======================================================================
% Fonte: analysis_summary_extreme_events_t1_t6.txt ed extreme_event_metrics_
% t_plus_1/6.csv della run a sei orizzonti (letti l'11 ottobre 2026); metodo nel
% notebook pvgis_sde_extreme_events.ipynb. Episodi: 23 aprile 05:10 -- 26 aprile
% 18:10 (86 ore, picco 4169.1 il 24 aprile 08:10, 559 finestre estreme);
% 28 giugno 04:10 -- 29 giugno 18:10 (39 ore, picco 5094.6 il 28 giugno 17:10,
% 500 finestre estreme). Riferimento normale: 4390625 righe per orizzonte.
% Classifica: aprile t+1 MAE 70.47 contro 42.62 (+65.3%), PICP 0.9464 e 0.9433,
% CLC 0.927; t+6 145.71 contro 75.42 (+93.2%), 0.9453 e 0.9314, CLC 1.510;
% giugno t+1 35.55 contro 45.62 (-22.1%), 0.9454 e 0.9632, CLC 0.517; t+6 58.62
% contro 91.01 (-35.6%), 0.9198 e 0.9706, CLC 0.889.
% Segno del bias: previsto meno osservato (nella fascia 0--20% delle normali
% bias +25.8 con quota di sovrastime 0.587). Aprile, t+1: fascia 60--80% bias
% -141.2, PICP 0.757; fascia 80--100% bias -138.5, PICP 0.869.
% DA VERIFICARE nel notebook sul server: quantile che definisce la coda e
% distanza massima fra ore dello stesso episodio.
% Nei file i due episodi si chiamano april_dust e june_extreme. Nel testo NON è
% indicata la causa meteorologica (polvere sahariana, ondata di calore): serve
% una fonte verificata prima di scriverla.
% Non usato: riepilogo_arrivo_evento (bias di picco 49.0 W a t+1 e 188.6 W a
% t+6, ritardi 0 e 4 ore, assestamento dopo 4 e 7 ore): non è indicato a quale
% episodio si riferisce.

% ======================================================================
% Errore in funzione della soglia e bande di dispersione
% ======================================================================
% Fonte: detector_threshold_sensitivity_metrics.csv (letto l'11 ottobre 2026).
% MTGFlow, colonne: k, normali, segnalate, quota, MAE normali, MAE segnalate,
% RMSE normali, RMSE segnalate, rapporto MAE.
%   t+1  0.5: 3972525, 752179, 15.92%, 44.448, 50.205, 68.864, 78.593, 1.1295
%   t+1  1.0: 4260191, 464513,  9.83%, 44.760, 50.910, 69.429, 79.686, 1.1374
%   t+1  1.5: 4412660, 312044,  6.60%, 44.965, 51.013, 69.806, 79.706, 1.1345
%   t+1  2.0: 4503292, 221412,  4.69%, 45.090, 50.941, 70.042, 79.296, 1.1298
%   t+1  2.5: 4558534, 166170,  3.52%, 45.161, 50.939, 70.175, 78.972, 1.1279
%   t+1  3.0: 4594272, 130432,  2.76%, 45.192, 51.425, 70.235, 79.385, 1.1379
%   t+6  0.5: 81.690, 87.269, 114.106, 123.987, 1.0683
%   t+6  1.0: 81.941, 88.426, 114.655, 125.210, 1.0792
%   t+6  1.5: 82.124, 89.004, 115.054, 124.970, 1.0838
%   t+6  2.0: 82.261, 89.034, 115.323, 123.835, 1.0823
%   t+6  2.5: 82.340, 89.118, 115.467, 122.876, 1.0823
%   t+6  3.0: 82.369, 89.959, 115.513, 123.309, 1.0922
% STGAN, colonne: quota del ranking, normali, segnalate, quota fra le diurne,
% MAE normali, MAE segnalate, RMSE normali, RMSE segnalate, rapporto MAE.
%   t+1 0.25: 4700621,  24083, 0.51%, 45.205, 76.532, 70.162, 119.689, 1.6930
%   t+1 0.50: 4676922,  47782, 1.01%, 45.128, 68.507, 70.024, 107.497, 1.5181
%   t+1 0.75: 4653385,  71319, 1.51%, 45.076, 64.154, 69.933, 100.956, 1.4232
%   t+1 1.00: 4630172,  94532, 2.00%, 45.041, 61.193, 69.873,  96.447, 1.3586
%   t+1 1.25: 4607233, 117471, 2.49%, 45.016, 59.031, 69.827,  93.220, 1.3113
%   t+1 1.50: 4584753, 139951, 2.96%, 45.000, 57.311, 69.797,  90.639, 1.2736
%   t+1 2.00: 4541473, 183231, 3.88%, 44.984, 54.783, 69.764,  86.844, 1.2178
%   t+1 3.00: 4465207, 259497, 5.49%, 45.014, 51.399, 69.790,  81.811, 1.1418
%   t+1 4.00: 4406476, 318228, 6.74%, 45.062, 49.544, 69.844,  79.072, 1.0994
%   t+1 5.00: 4362749, 361955, 7.66%, 45.115, 48.368, 69.902,  77.378, 1.0721
%   t+6 0.25: 82.479, 102.020, 115.591, 141.172, 1.2369
%   t+6 0.50: 82.427,  97.387, 115.518, 135.372, 1.1815
%   t+6 0.75: 82.382,  95.391, 115.453, 132.865, 1.1579
%   t+6 1.00: 82.349,  93.805, 115.409, 130.718, 1.1391
%   t+6 1.25: 82.318,  92.793, 115.369, 129.284, 1.1273
%   t+6 1.50: 82.289,  92.047, 115.337, 128.107, 1.1186
%   t+6 2.00: 82.221,  91.429, 115.264, 126.861, 1.1120
%   t+6 3.00: 82.116,  90.529, 115.178, 124.943, 1.1025
%   t+6 4.00: 82.029,  90.193, 115.111, 124.060, 1.0995
%   t+6 5.00: 81.959,  90.039, 115.064, 123.538, 1.0986
% analysis_metadata.json: sola rielaborazione, nessun nuovo addestramento;
% soglia diurna 10 W/m2; regola MTGFlow score >= Q3 + k IQR, riferimento k 1.5;
% regola STGAN top-K globale, riferimento 1%, ranking ricalcolato dopo
% l'esclusione di 18 istanti (isolated_regional_solar_dropout_plus_immediate_
% recovery); dominio comune MTGFlow/STGAN, copertura 0.9952; picco di
% riferimento per le fasce di produzione 712.95 W; tipi meteorologici dalle
% etichette climatologiche 2005--2018.
% Le righe STGAN sopra sono quelle della cartella anomaly_threshold_sensitivity_
% t1_t6, i cui dati di qualità vengono da pvgis_stgan/prepared: versione su
% grafo (DA CONFERMARE nel notebook).
% Variante ConvGRU, kernel 3 (anomaly_threshold_sensitivity_cnn/
% convgru_reference_seed_20_t1_t6_20260916T133237162286Z_45c3ff1b, letto
% l'11 ottobre 2026). Colonne: quota del ranking, segnalate, quota fra le
% diurne, MAE normali, MAE segnalate, rapporto.
%   t+1 0.25:  24147, 0.51%, 45.324, 53.248, 1.1748
%   t+1 0.50:  47814, 1.01%, 45.319, 49.836, 1.0997
%   t+1 0.75:  70798, 1.50%, 45.328, 47.776, 1.0540
%   t+1 1.00:  93388, 1.98%, 45.342, 46.445, 1.0243
%   t+1 1.25: 115766, 2.45%, 45.363, 45.419, 1.0012
%   t+1 1.50: 138081, 2.92%, 45.386, 44.655, 0.9839
%   t+1 2.00: 182487, 3.86%, 45.435, 43.614, 0.9599
%   t+1 3.00: 270849, 5.73%, 45.560, 42.152, 0.9252
%   t+1 4.00: 358113, 7.58%, 45.713, 41.108, 0.8993
%   t+1 5.00: 443468, 9.39%, 45.895, 40.240, 0.8768
%   t+6 0.25: 82.586, 81.029, 0.9811
%   t+6 0.50: 82.627, 77.827, 0.9419
%   t+6 0.75: 82.680, 75.905, 0.9181
%   t+6 1.00: 82.740, 74.577, 0.9013
%   t+6 1.25: 82.801, 73.706, 0.8902
%   t+6 1.50: 82.861, 73.178, 0.8831
%   t+6 2.00: 82.988, 72.381, 0.8722
%   t+6 3.00: 83.232, 71.825, 0.8629
%   t+6 4.00: 83.460, 71.830, 0.8607
%   t+6 5.00: 83.661, 72.125, 0.8621
% Non lette: righe MTGFlow e metadati di questa seconda analisi (dominio e
% previsioni usate): DA LEGGERE per confermare che il previsore è lo stesso.
% DA FARE: fasce di produzione e bande (quartili e baffi) degli errori assoluti,
% che richiedono una figura.
% TESTO SPOSTATO DAL CAPITOLO 5 (11 ottobre 2026). Descrive come è calcolata
% l'analisi: da riusare, accorciato, in apertura della sottosezione e soltanto
% insieme al risultato corrispondente.
% Nell'analisi di sensibilità MTGFlow, gli anomali soddisfano la soglia per sito $Q_3+k\,\mathrm{IQR}$, usando le statistiche di riferimento del training. Per STGAN si varia la percentuale del ranking globale. A ogni soglia vengono ricalcolati i gruppi e i rispettivi MAE e RMSE: l'asse orizzontale descrive la soglia, non il tempo.
%
% Le curve MAE sono accompagnate da due fasce degli errori assoluti: una fascia più scura fra primo e terzo quartile e una più chiara fra i baffi di Tukey. I baffi corrispondono agli estremi osservati entro $Q_1-1.5\,\mathrm{IQR}$ e $Q_3+1.5\,\mathrm{IQR}$; non comprendono necessariamente tutto l'intervallo dei valori. Gli outlier continuano a contribuire al MAE.
%
% Queste fasce descrivono la dispersione degli errori individuali. Non sono intervalli di confidenza del MAE e non sono intervalli predittivi della rete. La stessa distinzione vale per i boxplot degli errori assoluti nei quattro casi; il RMSE è rappresentato come metrica aggregata. La sovrapposizione delle fasce non costituisce un test di significatività.
%
% La valutazione probabilistica considera separatamente copertura e ampiezza degli intervalli predittivi del previsore. Il confronto fra STGAN su grafo e variante ConvGRU sulle coordinate comuni riporta inoltre concordanze, discordanze e Jaccard delle anomalie, distinguendo patch complete e incomplete. L'accordo fra detector non equivale ad accuratezza rispetto a una verità di riferimento indipendente.
% Inserire sensibilità MTGFlow e STGAN con curve MAE, fascia Q1--Q3 e baffi.
% Le bande descrivono errori individuali: nessuna conclusione di significatività
% o equivalenza dei gruppi può essere ricavata dalla sola sovrapposizione.

% ======================================================================
% Errori per tipo e direzione dell'estremo meteorologico
% ======================================================================
% Fonte: detector_threshold_sensitivity_by_meteo_type.csv, righe mtgflow k 1.5
% e stgan 1.0, orizzonti 1 e 6 (lette l'11 ottobre 2026). Colonne: numero,
% quota delle segnalate, MAE e RMSE a t+1, MAE e RMSE a t+6.
% MTGFlow (312044 segnalate su 4724704):
%   irradiance_high    2762, 0.89%,  39.66,  59.40; 124.82, 148.76
%   irradiance_low     6842, 2.19%, 117.84, 154.30; 231.53, 276.46
%   temperature_high  48695, 15.61%, 33.47,  59.17;  56.94,  82.16
%   temperature_low    1829, 0.59%,  69.95, 100.53; 147.78, 202.68
%   wind_high         20251, 6.49%,  68.29, 102.23; 115.30, 152.00
%   wind_low           5948, 1.91%,  55.94,  91.57;  93.27, 143.76
%   untyped          231716, 74.26%, 51.09,  77.83;  88.71, 122.25
% STGAN (94532 segnalate):
%   irradiance_high     167, 0.18%,  62.17, 110.69; 147.40, 170.05
%   irradiance_low     5571, 5.89%, 110.12, 169.28; 184.57, 231.84
%   temperature_high   2977, 3.15%,  40.31,  68.13;  63.07,  88.62
%   temperature_low    2356, 2.49%,  64.53,  90.41; 157.95, 196.59
%   wind_high          3700, 3.91%,  69.58,  99.41; 123.31, 161.33
%   wind_low           1594, 1.69%,  72.75, 120.75;  94.01, 132.82
%   untyped           79668, 84.28%, 58.13,  89.90;  86.72, 119.71
% Tipi dalle etichette climatologiche 2005--2018 (analysis_metadata.json):
% alto = valore sopra climatology_q_high, basso = sotto climatology_q_low.
% Le righe con poche osservazioni (irradianza alta per STGAN: 167) vanno lette
% con cautela.
% TESTO SPOSTATO DAL CAPITOLO 5 (11 ottobre 2026). Descrive come è calcolata
% l'analisi: da riusare, accorciato, in apertura della sottosezione e soltanto
% insieme al risultato corrispondente.
% Le anomalie STGAN possono essere ulteriormente suddivise intersecando la decisione del detector con flag climatologici delle variabili meteorologiche, alla stessa località e al timestamp target. Per temperatura e vento vengono distinti i lati alto e basso confrontando il valore con le bande climatologiche della relativa osservazione; per l'irradianza viene verificata anche la variabile associata al flag. Il criterio non confonde quindi un estremo della produzione con un estremo dell'irradianza.
%
% I sottogruppi possono sovrapporsi. Una stessa osservazione può soddisfare più condizioni meteorologiche, mentre un'anomalia non associata ad alcuna condizione rimane anomala e viene conteggiata separatamente. I normali conservano la decisione STGAN e non vengono ridefiniti come complemento di ciascuna categoria. Le statistiche per categoria non costituiscono pertanto una partizione disgiunta di tutte le osservazioni.
%
% L'analisi degli errori tipizzati considera i target diurni con irradianza maggiore di $10\,\mathrm{W\,m^{-2}}$. Per le fasce di produzione utilizza un riferimento condiviso fra orizzonti: il percentile 99 della produzione positiva diurna, contando una sola volta ciascuna coppia località--timestamp target, e intervalli di ampiezza 20 punti percentuali. Oltre agli errori aggregati, calcola il RMSE per località e ne riassume mediana e quartili. La banda fra siti descrive l'eterogeneità geografica del RMSE; è distinta dai quartili degli errori assoluti individuali e dagli intervalli di confidenza di una metrica.
% Riportare normali STGAN ed estremi alti/bassi di temperatura, vento e
% irradianza, ottenuti intersecando anomalie e flag climatologici al target.
% Esplicitare filtro diurno, riferimento di produzione e orizzonti confrontati.
% Contare sottogruppi sovrapposti e anomalie non tipizzate senza trasformare
% queste ultime in normali. La somma dei gruppi non è il totale dei campioni.
% Distinguere MAE/RMSE aggregati, quartili degli errori assoluti individuali
% e mediana/quartili del RMSE fra località; nessuna banda è automaticamente
% un intervallo di confidenza.

% ======================================================================
% Costo computazionale e complessità dei modelli
% ======================================================================
% Inserire parametri e tempi misurati su hardware e carichi confrontabili.
% Separare preprocessing, training e scoring. La ConvGRU conserva la LSTM lunga;
% la rimozione della GAT dal previsore non rimuove la costruzione del grafo.
% Distinguere effetti di kernel e numero di parametri da quelli di batch,
% cache e memmap; non attribuire accelerazioni senza misure confrontabili.
% La modalità block di shuffle conserva i campioni ma ne cambia l'ordine.
% Per ERA5 distinguere numero di aggiornamenti per epoca e costo per cella:
% nelle varianti sul campo intero e su grafo il discriminatore resta a patch,
% quindi la riduzione degli aggiornamenti non implica una riduzione
% proporzionale del tempo. Riportare tempi e memoria misurati sulle run reali,
% con precisione numerica e hardware indicati.
% Limiti e discussione complessiva: Capitolo delle conclusioni.
